\documentclass[11pt]{article}
\usepackage[margin=1in]{geometry}
\usepackage{xcolor}
\usepackage{hyperref}
\usepackage{enumitem}
\usepackage{soul}

% Reviewer comments shown in indented italics
\newenvironment{reviewer}{%
  \begin{quote}\itshape\color{black!70}}{%
  \end{quote}}

% Author responses shown plain
\newcommand{\response}[1]{\noindent\textbf{Response:} #1\par\medskip}

% Manuscript changes shown bulleted
\newenvironment{changes}{%
  \noindent\textbf{Changes to manuscript:}
  \begin{itemize}[leftmargin=2em,itemsep=2pt]}{%
  \end{itemize}\medskip}

\title{Response to Reviewers\\
\large Manuscript ID: ci-2026-009433\\
``When Do Simple Models Win? Machine Learning Architectures for UV Absorption Prediction''}
\author{Umesh Arampath, Bryant Pero, David Stewart, David Demirjian}
\date{\today}

\begin{document}
\maketitle

We thank the Associate Editor and both reviewers for their thoughtful and constructive feedback. The reviewers' comments have substantially improved the manuscript, particularly by sharpening the scope of our claims, strengthening the data preprocessing and evaluation protocol, and adding analyses that contextualize our findings.

In the responses below, each reviewer comment is reproduced verbatim in indented italics, followed by our response and a bulleted list of corresponding changes to the manuscript. All page and line numbers refer to the revised manuscript unless stated otherwise.

\section*{Summary of Major Changes}
\begin{itemize}[leftmargin=2em]
  \item Retitled the paper and revised the abstract and conclusion to restrict generality claims to UV absorption and closely related local optical properties; removed the ``Insights for General Molecular Property Prediction'' framing.
  \item Removed all ``RF is optimal'' language; reframed the RF/D-MPNN comparison as a statistical tie ($p = 0.77$), with RF distinguished by compute cost and operational simplicity rather than predictive accuracy.
  \item Added a new Results subsection (``Performance on Non-Local Subsets'') reporting per-model error metrics on $\lambda_{\max} > 500$~nm, $> 600$~nm, donor--acceptor, D-A-D, azo, and BODIPY subsets, with full SMARTS rules, subset counts, and a complete RMSE/MAE/$R^2$ breakdown added to the Supporting Information.
  \item Added Greenman et al.\ (\textit{Chem.\ Sci.} \textbf{2022}, \textit{13}, 1152) to Related Work and corrected the manuscript's previous overclaim that systematic solvent-encoding evaluations for UV absorption had been ``limited.''
  \item Quantified the chemical-space overlap between the primary Joung+Beard dataset and the Deep4Chem ``external'' benchmark (77.4\% of D4C solutes / 72.6\% of (solute, solvent) pairs are also in J+B) and reframed Deep4Chem as in-distribution rather than OOD validation, with the more rigorous Jung 2024 benchmark ($\sim$50\% solute overlap) carrying the OOD claim. Added a new SI section with the full overlap analysis and Venn diagram.
  % TODO: fill in as we work through the remaining comments
\end{itemize}

\clearpage

\section*{Reviewer 1}

\subsection*{Comment R1-1: Generality of conclusions}
\begin{reviewer}
The authors make some claims about the conclusions from this work being extensible to molecular property prediction in general (e.g.\ in the title, pg.\ 1 line 35 in abstract, pg.\ 39 line 34 in conclusion). While this work seems to give evidence to the authors' hypothesis regarding the locality of target properties, I think these statements go too far given that the work does not validate this hypothesis in any domains other than optical properties. The work frames UV absorption as ``an ideal benchmarking target'' (pg.\ 3 line 25), and while the stated reasons do indeed make UV absorption a particularly interesting benchmarking task, I'm not sure that it's appropriate to say that it's ``ideal'' in the sense that conclusions that we draw from benchmarking on UV absorption can necessarily be applied more generally in molecular property prediction.
\end{reviewer}

\response{We agree with the reviewer that our generality claims were not supported by the experiments reported in this manuscript, which evaluate only UV absorption $\lambda_{\max}$. We have removed all references to ``general molecular property prediction'' from the title, abstract, introduction, and conclusion, softened the ``ideal benchmarking target'' framing, and recast the locality argument as a testable hypothesis whose cross-property validation is left as future work.}

\begin{changes}
  \item \textbf{Title.} Removed the trailing clause ``and Insights for General Molecular Property Prediction.'' The title now reads: ``When Do Simple Models Win? Machine Learning Architectures for UV Absorption Prediction.''
  \item \textbf{Abstract (closing paragraph).} Replaced the paragraph beginning ``These findings provide general guidance for molecular property prediction\ldots'' with a UV-scoped restatement and an explicit hypothesis sentence: ``Within UV absorption prediction, optimal model choice depends on the relationship between target compounds and available training data: simple fingerprint-based models are an efficient and accurate alternative to deep learning for interpolation within known chemical space, while learned representations from graphs, sequences, or pretraining generalize better to novel scaffolds. We further hypothesize that this pattern\ldots\ may extend to other locally determined molecular properties; validating this across property types is left as future work.''
  \item \textbf{Introduction (``ideal benchmarking target'').} Softened the framing from ``The UV absorption prediction task presents several distinctive challenges that make it an ideal benchmarking target'' to ``The UV absorption prediction task is well-suited for systematic architecture comparison for several reasons.''
  \item \textbf{Conclusion (``More broadly'' paragraph).} Removed the three sentences claiming generality to molecular property prediction beyond UV absorption and replaced them with: ``Our results raise a testable hypothesis about model selection: that architectures whose inductive bias matches the locality of the target property will outperform global-attention architectures in-distribution, while pretrained representations confer out-of-distribution advantages regardless of inductive bias. The present work establishes this pattern for UV absorption $\lambda_{\max}$; whether it extends to fluorescence emission, lipophilicity, aqueous solubility, or other molecular endpoints whose locality differs from UV absorption requires dedicated cross-property benchmarking, which we identify as a priority for future work.''
  \item \textbf{Supporting Information title.} Updated the SI title to match the revised main-manuscript title (``Supporting Information: When Do Simple Models Win? Machine Learning Architectures for UV Absorption Prediction''), dropping the trailing ``and Insights for General Molecular Property Prediction'' clause.
  \item The existing Future Work section already frames cross-property extension as a hypothesis to validate; it is retained without further change.
\end{changes}

\subsection*{Comment R1-2: ``RF is optimal'' is hurt by the absence of a statistically significant difference between RF and Chemprop}
\begin{reviewer}
The claim that RF is ``optimal'' (pg 1 line 51 in abstract) is hurt by the fact that the difference between RF and Chemprop is not statistically significant.
\end{reviewer}

\response{We agree. We have removed every instance of ``RF is optimal'' or equivalent language and replaced it with statements that explicitly acknowledge the statistical equivalence between RF and the D-MPNN ($p = 0.77$). The revised text frames RF's advantage as one of compute cost and operational simplicity rather than predictive accuracy, which our data support.}

\begin{changes}
  \item \textbf{Abstract (line 53).} Replaced ``RF with Morgan fingerprints is optimal: it matches the GNN\ldots'' with ``RF with Morgan fingerprints is competitive with the best deep learning model at a fraction of the compute cost: RF and the D-MPNN are statistically indistinguishable (RMSE = $31.34 \pm 1.82$ vs.\ $31.69 \pm 3.18$~nm, $p = 0.77$), and RF trains in 15 minutes on a CPU with no hyperparameter tuning while outperforming all sequence-based models.''
  \item \textbf{Introduction contributions bullet.} Replaced ``RF is optimal for in-distribution screening (fast, no GPU, no tuning needed)\ldots'' with ``on in-distribution screening, RF is statistically tied with the best deep learning model (D-MPNN, $p = 0.77$) but trains $\sim$100$\times$ faster with no hyperparameter tuning; on novel compounds, deep learning models generalize better.''
  \item \textbf{Figure caption (Training time vs.\ performance).} Replaced ``RF achieves the best RMSE in 15~minutes\ldots'' with ``RF attains performance statistically tied with the best deep learning model in 15~minutes\ldots,'' eliminating the implicit optimality claim while preserving the compute-cost comparison the figure is designed to convey.
\end{changes}

\subsection*{Comment R1-3: ``3D methods offering advantages'' typo}
\begin{reviewer}
Pg.\ 3 line 39: ``3D methods offering advantages primarily for structurally diverse datasets (see Cross-Dataset Benchmarking)'' --- Is this a typo referring to 3D methods? This work doesn't use 3D methods, as discussed in the limitations and future work sections.
\end{reviewer}

\response{The reviewer is correct. The phrase was a leftover from an earlier draft and is inaccurate. We have removed the ``3D methods offering advantages\ldots (see Cross-Dataset Benchmarking)'' clause, corrected ``1D representations'' to ``1D and 2D representations'' to accurately describe the architectures benchmarked (which include the D-MPNN operating on molecular graphs), and explicitly state that 3D methods are not evaluated in this work, with a pointer to the Limitations and Future Work sections.}

\begin{changes}
  \item \textbf{Introduction bullet (Representation diversity).} The bullet now reads: ``Although absorption fundamentally depends on electronic structure influenced by molecular geometry, our results show that 1D and 2D representations (Morgan fingerprints, SMILES strings, molecular graphs) capture the dominant structure--property relationships for UV absorption without requiring explicit 3D coordinates. Methods that incorporate 3D conformational information are not evaluated in this work and are discussed in the Limitations and Future Work sections.''
\end{changes}

\subsection*{Comment R1-4: ``Architecture-independent solvent encoding'' mischaracterization}
\begin{reviewer}
Pg.\ 4 line 33: I'm not sure if it's appropriate to say ``Architecture-independent solvent encoding'' when some architectures use Morgan fingerprints to represent solvents and others learn the representation. If I understand correctly, the only thing that's architecture independent is that the solvent representation is simply concatenated with the solute representation rather than using some more complicated process to combine them.
\end{reviewer}

\response{The reviewer is correct. Each architecture uses its native representation for both solute and solvent (Morgan fingerprints for the tree ensembles, character-level SMILES for the sequence models, and separate D-MPNN graph encoders for Chemprop). What is consistent across architectures is the concatenation \emph{strategy}, not the encoding itself. We have rephrased the contribution accordingly. The detailed Solvent Information Encoding section already describes each architecture's native encoding scheme, so no change is needed there.}

\begin{changes}
  \item \textbf{Introduction contributions bullet.} Replaced ``Architecture-independent solvent encoding: a simple concatenation strategy yields 6--8\% RMSE reductions\ldots'' with ``Concatenation-based solvent encoding generalizes across architectures: the same strategy of concatenating native solute and solvent representations---Morgan fingerprints for the tree ensembles, SMILES strings for the sequence models, and molecular graphs for the D-MPNN---yields 6--8\% RMSE reductions across all three architecture families ($p < 0.02$), without domain-specific descriptor engineering.''
\end{changes}

\subsection*{Comment R1-5: Non-local compounds (D-A, D-A-D, long wavelength) subset analysis}
\begin{reviewer}
The idea that absorption is a locally determined property is right to a first approximation, but in some cases (e.g.\ longer-wavelength dyes, donor-acceptor or donor-acceptor-donor dyes, etc.) global effects, charge transfer, and 3D geometry become more important. It would be insightful if the authors could identify any compounds in their dataset where non-local effects would be expected to be important, and see if these results lend more support to their claims about locality and architecture choice.
\end{reviewer}

\response{We thank the reviewer for this suggestion---it is exactly the kind of test that can either strengthen or honestly weaken the locality argument. We have added a new Results subsection (``Performance on Non-Local Subsets,'' immediately preceding the Error Analysis subsection) that partitions the Joung+Beard test folds into subsets where non-local effects are expected to dominate and reports per-model error metrics on each subset. The subsets are: long-wavelength compounds ($\lambda_{\max} > 500$~nm and $> 600$~nm), donor--acceptor (D-A) molecules, donor--acceptor--donor (D-A-D) molecules, azo dyes (Ar--N$=$N--Ar), and BODIPYs (BF$_2$ bridging two ring nitrogens). Donors and acceptors are identified via 7 and 12 SMARTS rules respectively (full patterns in Supporting Information Tables~S10--S12, with the acceptor rule generalised to match aromatic-or-alkene $sp^2$ anchors so that push--pull cinnamic acids such as ferulic acid are captured). The union of all non-local categories covers 49\% of the primary dataset (9{,}229 / 18{,}755 molecules).

The results give mixed support for the locality argument. On the full dataset, RF and Chemprop are statistically indistinguishable (RMSE = 31.3 vs.\ 31.7~nm, $p = 0.77$), as previously reported. On every subset enriched for non-local effects, Chemprop edges out RF: by 1.2~nm on D-A, 2.4~nm on D-A-D, 3.4~nm on azo, and 6.7~nm on $\lambda > 500$~nm; on $\lambda > 600$~nm the gap is 8.6~nm (Chemprop RMSE = 60.9 vs.\ RF 69.5, both with large per-fold variance). This is the direction predicted by the locality hypothesis---a message-passing GNN should benefit more, relative to a fingerprint-based model, when the property depends on longer-range electronic effects. However, the absolute gaps remain modest and within one fold-standard-deviation for the smaller subsets, so we present this as suggestive rather than conclusive evidence and have phrased the discussion that way in the manuscript. The one row where RF holds its ground is BODIPYs, which makes physical sense: the BODIPY scaffold is rigid and conserved, so absorption is well-predicted by short-range substitution patterns that fingerprints already capture.

In the MAE column (Supporting Information Table~S14), the picture is more nuanced: Chemprop only wins on MAE for the long-wavelength and azo subsets, while RF has lower MAE on D-A, D-A-D, BODIPYs, stilbenes, and cyanines. This reflects RF's tendency to produce residual distributions concentrated near zero with a longer tail of larger errors, which is penalised more under squared-error than under absolute-error. The reviewer's intuition is therefore borne out specifically for the squared-error metric on long-wavelength chromophores.

The first trend---all five architectures degrading sharply on $\lambda > 500$~nm---is consistent with prior observations by Beard et al.\ (2019) and Joung et al.\ (2020--2021) on the same datasets, and stems from data scarcity at long wavelengths combined with the greater structural complexity of low-gap chromophores (extended polyenes, push--pull cyanines, BODIPYs, phthalocyanines). ChemBERTa exhibits an extreme version of this failure mode: its predictions collapse to an essentially constant value around 600~nm, giving negative $R^2$ on both long-wavelength subsets. This is the same prediction ceiling visible in Figure~5 (parity plot) and addressed separately in our response to R1-16.

All numbers in the new subsection are reproducible from the per-fold predictions distributed with the code via\\
\hspace*{1em}\texttt{python3 analyze\_nonlocal\_subsets.py --models rf\_tuned xgboost\_v2 chemprop\_v2 bigru\_tuned chemberta --regen\_xgboost}\\
which writes the underlying CSVs to \texttt{results/nonlocal\_subsets/}.}

\begin{changes}
  \item \textbf{Results---new subsection.} Added ``Performance on Non-Local Subsets'' as a new subsection of Results and Discussion, immediately preceding ``Error Analysis.'' The subsection defines five non-local subsets via SMARTS, reports an in-text RMSE table (Table~5 in the revised manuscript), and presents the two-trend discussion summarised above.
  \item \textbf{Supporting Information---new section.} Added ``SMARTS Rules and Metrics for Non-Local Subset Analysis'' (Section~S7) containing five tables: the donor SMARTS list (Table~S10), the acceptor SMARTS list (Table~S11), the dye-motif SMARTS list (Table~S12), the full-dataset subset counts (Table~S13), and a complete RMSE/MAE/$R^2$ breakdown across all subsets and all five models (Table~S14).
  \item \textbf{Code release.} Added \texttt{analyze\_nonlocal\_subsets.py} to the repository: it loads the saved per-fold predictions for the four DL/ensemble models, re-scores the XGBoost models from their joblib files, applies the SMARTS classification, and writes the per-fold and aggregated CSVs that the manuscript tables are built from.
\end{changes}

\subsection*{Comment R1-6: Cite prior solvent-encoding benchmark (Greenman et al., 2022)}
\begin{reviewer}
Page 8 line 41: ``but systematic evaluation of solvent encoding strategies for UV absorption prediction has been limited.'' --- Prior work benchmarked 4 different solvent encoding strategies for UVVis absorption (\url{https://doi.org/10.1039/d1sc05677h})
\end{reviewer}

\response{We thank the reviewer for pointing out this oversight; the cited paper is Greenman, Green, and G\'omez-Bombarelli, ``Multi-fidelity prediction of molecular optical peaks with deep learning,'' \textit{Chem.\ Sci.} \textbf{2022}, \textit{13}, 1152--1162 (DOI 10.1039/D1SC05677H). It is the most directly relevant prior benchmark of solvent-encoding strategies for UV/Vis absorption and our previous text understated its contribution. We have added a Related Work paragraph that names the four strategies Greenman et al.\ compared (no solvent; one-hot solvent labels; Morgan-fingerprint concatenation; a second D-MPNN encoder for the solvent), and rephrased the contribution claim so that we no longer assert that systematic evaluation has been ``limited.'' Instead, we now state that Greenman et al.\ established that explicit solvent representation matters \emph{within a D-MPNN architecture} for UV/Vis peak prediction, and that the open question we address here is the comparative behaviour of fingerprint- and sequence-based architectures under the same concatenation-based scheme on a shared dataset. The Greenman et al.\ reference is also cited later in our response to R1-8 (duplicate handling) and R1-9 (chromophore-solvent leakage), where their methodology is directly relevant.}

\begin{changes}
  \item \textbf{Related Work (Solvent effects, pg.\ 8 line 41).} Replaced ``but systematic evaluation of solvent encoding strategies for UV absorption prediction has been limited'' with a new two-sentence anchor that introduces Greenman et al.~\cite{greenman2022multifidelity}, names the four solvent-encoding strategies they benchmarked, and reframes our contribution as a cross-architecture rather than novel-strategy comparison.
  \item \textbf{Bibliography.} Added \texttt{greenman2022multifidelity} (Greenman, Green, G\'omez-Bombarelli, \textit{Chem.\ Sci.} \textbf{2022}, \textit{13}, 1152) to \texttt{Proposal.bib}. Also added \texttt{song2025fluortools} (Song et al., \textit{Digital Discovery} \textbf{2025}, \textit{4}, 3728) in anticipation of our responses to R1-7 and R1-8, where their Table 3 of UV absorption datasets and their duplicate-handling protocol are directly relevant.
\end{changes}

\subsection*{Comment R1-7: Deep4Chem overlap with Joung; choice of Joung+Beard as primary dataset}
\begin{reviewer}
The Deep4Chem dataset isn't a great ``external dataset'' given that it substantially overlaps with the Joung dataset that makes up part of the primary dataset. Why were the Joung and Beard datasets chosen to make the primary dataset as opposed to the many other UV absorption datasets that exist (see Table 3 here: \url{https://doi.org/10.1039/D5DD00402K})?
\end{reviewer}

\response{The reviewer makes two related points, both of which we accept and address with new analysis and text.

\textbf{(a) On the Deep4Chem overlap.} The reviewer is correct: Deep4Chem was the source dataset for the Joung half of our primary collection, so by construction it shares a very high fraction of records with Joung+Beard. To replace hand-waving with a number, we canonicalised the SMILES of both datasets through RDKit and computed set-intersections at three levels of granularity. The results, now reported in Supporting Information Section~S8 (Table~S15 and Figure~S3), are:
\begin{itemize}[leftmargin=2em,itemsep=2pt]
  \item \emph{Unique canonical solute SMILES:} 5{,}018 of Deep4Chem's 6{,}480 unique solutes (77.4\%) also appear in Joung+Beard. Conversely, 63.8\% of Joung+Beard solutes appear in Deep4Chem.
  \item \emph{Unique (solute, solvent) pairs (both canonicalised):} 12{,}654 of Deep4Chem's 17{,}427 pairs (72.6\%) match a pair in Joung+Beard exactly.
  \item \emph{Generic Bemis--Murcko scaffolds:} 2{,}467 of Deep4Chem's 2{,}498 scaffolds (98.8\%) are also present in Joung+Beard; the two datasets explore essentially the same scaffold space.
\end{itemize}
These numbers vindicate the reviewer's intuition that Deep4Chem is poorly characterised as an OOD external benchmark. We have updated the Methods (Dataset Description and Preparation) to state explicitly that ``Deep4Chem is therefore best characterised as an \emph{in-distribution external} benchmark---useful for reproducing published per-paper baselines under their own splits, but \emph{not} a test of out-of-distribution generalisation.'' We have also added a sentence to the Cross-Dataset Benchmarking subsection noting that the Deep4Chem results should be read as in-distribution validation, and that the Jung 2024 results (which we also report) constitute the more rigorous external test because Jung 2024 has only 48.8\% solute overlap with Joung+Beard, a separate experimental provenance, and distinct chromophore families (notably long-wavelength dyes from the solar-cell and OLED literature).

\textbf{(b) On the choice of Joung+Beard as the primary dataset.} The reviewer rightly points out that Song et al.~\cite{song2025fluortools} survey nine UV/Vis dye datasets in their Table~3 (ChemFluor, Deep4Chem, SMFluo, ChemDataExtractor, Dye aggregation, DSSCDB, Fluor-predictor, Ksenofontov's data, Literature retrieval). We have added a four-part rationale to the Methods section explaining why Joung+Beard was selected as the primary dataset:
\begin{itemize}[leftmargin=2em,itemsep=2pt]
  \item It is the largest unified collection of \emph{purely experimental} solute--solvent $\lambda_{\max}$ measurements available without paywall or license restrictions at the time of submission, with each record carrying an experimental wavelength and an explicit solvent label.
  \item It has \emph{published baselines from independent groups} (Joung et al.\ GCNN on Deep4Chem, Beard et al.\ DR-CNN, Jung et al.\ multifidelity and GBFS) that enable apples-to-apples comparison without training every baseline ourselves.
  \item Combining the two source collections yields broader chromophore coverage than either alone (Deep4Chem 6{,}480 and ChemFluor 4{,}386 unique solutes vs.\ approximately 7{,}000 in the curated Joung+Beard primary set after our preprocessing).
  \item The dataset is dominated by drug-like organic chromophores rather than the dye-aggregation or solar-cell-specific compositions of DSSCDB and Dye-aggregation, matching our intended downstream applications (photosafety screening and wetlab UV-absorber discovery).
\end{itemize}
We considered but did not include Mamede's MEC dataset in the primary collection because of the Reaxys licence restrictions that prevent direct redistribution (see our response to R1-15); we use it instead as a downstream photosafety-classification benchmark where the labels and not the raw $\lambda_{\max}$ values are the supervision target. Computational datasets (e.g.\ Fluor-predictor's TD-DFT calculations) were excluded so that the primary regression target is purely experimental.

The script that produces the overlap numbers (\texttt{analyze\_deep4chem\_overlap.py}) is distributed with the code, so anyone can re-run the analysis or extend it to other dataset pairings.}

\begin{changes}
  \item \textbf{Methods (Dataset Description and Preparation, pg.\ 8).} Added a new four-part rationale paragraph immediately after the opening sentence, explaining why Joung+Beard was chosen as the primary dataset and citing Song et al.\ for the dataset-landscape survey.
  \item \textbf{Methods (Dataset Description and Preparation, pg.\ 8).} Extended the Deep4Chem description to report the measured overlap (77.4\% of solutes, 72.6\% of pairs) and to reframe Deep4Chem explicitly as an in-distribution rather than OOD external benchmark, with a pointer to the new SI overlap analysis (Section~S8, Figure~S3).
  \item \textbf{Methods (Dataset Description and Preparation, pg.\ 8).} Extended the Jung 2024 description to note the 48.8\% solute overlap and the qualitatively different chromophore coverage that makes it the more rigorous of the two external benchmarks.
  \item \textbf{Results (Cross-Dataset Benchmarking, pg.\ 30).} Added a caveat sentence at the end of the section noting that the Deep4Chem results reflect in-distribution interpolation rather than OOD generalisation, and that Jung 2024 carries the OOD claim.
  \item \textbf{Supporting Information---new section.} Added ``Overlap Between Primary and Deep4Chem Datasets'' (Section~S8) containing the per-granularity overlap table (Table~S15) and a side-by-side Venn diagram of solutes and (solute, solvent) pairs (Figure~S3); reproduced exactly by the supplied \texttt{analyze\_deep4chem\_overlap.py} script.
  \item \textbf{Figure renumbering housekeeping.} Insertion of the new Figure~S3 (Venn diagram) shifts the existing BiGRU character-type saliency figure from S3 to S4; the three main-text references to ``Figure~S3'' that pointed to the saliency analysis have been updated to ``Figure~S4'' accordingly.
\end{changes}

\end{document}
